\documentclass[preview]{standalone}

\usepackage{amsmath}
\usepackage{amssymb}
\usepackage{stellar}
\usepackage{definitions}

\begin{document}

\id{il-comportamentismo}
\genpage

\section{Il comportamentismo}

\begin{snippetdefinition}{comportamentismo-definition}{Comportamentismo}
    Il \emph{comportamentismo} è una scuola psicologica nata negli Stati Uniti
    che pone al centro della psicologia il comportamento osservabile. Rifiuta
    la tradizione mentalista e l'introspezione, sostenendo che la psicologia
    debba usare metodi controllabili come quelli delle scienze naturali.
\end{snippetdefinition}

\begin{snippetdefinition}{black-box-definition}{Black box}
    Per il comportamentismo, la mente è una \emph{scatola nera}, o
    \emph{black box}: lo psicologo non può né deve entrare direttamente al suo
    interno. L'indagine si concentra sulle relazioni osservabili tra stimoli
    ambientali e risposte comportamentali.
\end{snippetdefinition}

\begin{snippet}{modello-stimolo-risposta}
    L'obiettivo dello psicologo comportamentista è studiare associazioni
    \texttt{S-R}, dove \texttt{S} indica gli stimoli ambientali che impattano
    sulla black box e \texttt{R} indica le risposte dell'organismo. Il metodo
    sperimentale manipola la variabile indipendente, cioè lo stimolo, e misura
    la variabile dipendente, cioè la risposta comportamentale.
\end{snippet}

\begin{snippetexample}{thorndike-problem-box-example}{Problem box di Thorndike}
    Thorndike studia l'apprendimento inserendo gatti in una \emph{Problem box}.
    La gabbia contiene un meccanismo che permette all'animale di uscire. Il
    gatto, muovendosi casualmente, finisce per azionare il meccanismo. Con il
    ripetersi delle prove, il comportamento utile diventa più probabile. Questo
    esperimento illustra l'apprendimento per prove ed errori, definito
    apprendimento alla cieca perché l'animale non comprende inizialmente la
    soluzione.
\end{snippetexample}

\begin{snippetdefinition}{legge-effetto-definition}{Legge dell'effetto}
    Secondo la \emph{legge dell'effetto}, un animale tende a ripetere un'azione
    quando tale azione produce un effetto favorevole nell'ambiente.
\end{snippetdefinition}

\begin{snippetdefinition}{legge-esercizio-definition}{Legge dell'esercizio}
    La \emph{legge dell'esercizio} afferma che la ripetizione delle prove
    rafforza l'associazione tra una situazione e la risposta efficace.
\end{snippetdefinition}

\begin{snippetdefinition}{condizionamento-classico-definition}{Condizionamento classico}
    Il \emph{condizionamento classico}, proposto da Pavlov, è una forma di
    apprendimento in cui uno stimolo inizialmente neutro diventa capace di
    produrre una risposta dopo essere stato associato ripetutamente a uno
    stimolo incondizionato.
\end{snippetdefinition}

\begin{snippetexample}{pavlov-condizionamento-classico-storico-example}{Condizionamento classico di Pavlov}
    Pavlov osserva che il cane saliva non solo quando riceve cibo, ma anche
    quando segnali associati al cibo anticipano la somministrazione. Nella
    situazione iniziale, il cibo è lo stimolo incondizionato e produce
    salivazione, cioè una risposta incondizionata. Associando ripetutamente un
    suono al cibo, il suono diventa uno stimolo condizionato e finisce per
    produrre salivazione anche in assenza del cibo.
\end{snippetexample}

\includesnpt[width=75\%,src=/snippet/static-psychology/pavlov-classical-conditioning-overview.jpg]{centered-img}

\begin{snippetdefinition}{condizionamento-operante-definition}{Condizionamento operante}
    Il \emph{condizionamento operante}, elaborato da Skinner, è una forma di
    apprendimento in cui un comportamento aumenta o diminuisce in funzione
    delle conseguenze che produce nell'ambiente.
\end{snippetdefinition}

\begin{snippetdefinition}{rinforzo-definition}{Rinforzo}
    Il \emph{rinforzo} è uno stimolo che aumenta la probabilità che un certo
    comportamento si manifesti nuovamente.
\end{snippetdefinition}

\begin{snippet}{skinner-condizionamento-operante-storico}
    Skinner studia il condizionamento operante con la \emph{Skinner Box}, una
    gabbia in cui leve e rinforzi sono controllati sperimentalmente. Il
    soggetto può ricevere acqua, cibo o punizioni in funzione delle risposte
    emesse. Il metodo permette di studiare risposte complesse perché ciascuna
    risposta può essere seguita da un rinforzo.
\end{snippet}

\begin{snippetdefinition}{apprendimento-latente-definition}{Apprendimento latente}
    L'\emph{apprendimento latente} è una forma di apprendimento che avviene
    anche in assenza di rinforzo e che si manifesta solo successivamente,
    quando diventa utile per affrontare una situazione problematica.
\end{snippetdefinition}

\begin{snippet}{tolman-mappe-cognitive}
    Tolman sostiene che il rinforzo è utile perché il comportamento venga
    manifestato, ma non è necessario affinché venga appreso. Nei suoi studi sui
    labirinti, i topi costruiscono una rappresentazione spaziale dell'ambiente,
    detta mappa cognitiva, e la usano quando il percorso viene modificato.
\end{snippet}

\end{document}
